\honest{New Improved Lottery}{Eliezer Yudkowsky}{2007}

People are still saying that the lottery sells fantasy cheaply. One commenter wrote: ``There is a big difference between zero chance of becoming wealthy, and epsilon. Buying a ticket allows your dream of riches to bridge that gap.'' \nb{The comment continues: ``nobody can measure the value of entertainment or distant hope for another.'' The post does not quote that part.}

But between zero and epsilon there is ``an order-of-epsilon difference.'' If you doubt it, let epsilon be one over a googolplex. \nb{That is true of the probabilities. The commenter was talking about the dream, and whether a dream's worth should shrink with its probability is what the two disagree about.}

Then I design the New Improved Lottery. You pay one dollar once. The prize is paid about once every five years, at a random moment set by radioactive decay, so the phone could ring at any minute. When the boss asks you to rework a proposal, you can stare at the phone instead. You can fill a shopping cart and put everything back if no call comes. A screen on your cellphone could show the odds changing minute by minute. This dream would crowd out not only hopes of technical school but ``even hopes of getting home from work early.'' \nb{The harm pictured here is the previous post's claim, that lottery dreams displace better ones, made larger. That claim was offered there with ``maybe'' and ``might'' and no evidence, and the satire adds none.}

Then, in sarcasm: ``People are willing to pay; it must be valuable.''

Governments, with their ``vile monopoly on lotteries,'' do not sell this because they want you to pay every week. \nb{The United Kingdom has sold something close since 1956: Premium Bonds, which enter one purchase in every monthly prize draw until it is cashed, and whose price the government pays back on request.} So if the lottery is a service, ``it is clearly an enormously overpriced service,'' charged to the poorest, and it is your duty to demand the New Improved Lottery. \nb{On that premise the arithmetic agrees. In 2007 a ticket cost \$1 and each big game drew twice a week, so a ticket for every draw of one game came to about \$520 over five years.}
